\documentclass[10pt,a4paper]{amsart}
\usepackage[margin=19mm]{geometry}
\usepackage{amsmath,amssymb,mathtools}
\usepackage[scr=rsfs,cal=euler]{mathalfa}
\usepackage{xfrac}
\usepackage[unicode,hidelinks,bookmarks=false]{hyperref}
\newcommand{\ie}{i.e.\ }
\newcommand{\cf}{cf.\ }
\newcommand{\resp}{respectively\ }
\newcommand{\Subsection}{\subsection}
\providecommand{\nobreakdash}{\nobreak\hbox{-}}
\setlength{\emergencystretch}{2em}
\multlinegap0pt
\theoremstyle{plain}
\newtheorem{theorem}{Theorem}
\newtheorem{lemma}[theorem]{Lemma}
\newtheorem{proposition}[theorem]{Proposition}
\newtheorem{corollary}[theorem]{Corollary}
\theoremstyle{definition}
\newtheorem{definition}[theorem]{Definition}
\theoremstyle{remark}
\newtheorem{remark}[theorem]{Remark}
\newcommand{\RR}{\mathbb R}
\newcommand{\CC}{\mathbb C}
\newcommand{\BB}{\mathcal B}
\newcommand{\Bor}{\mathfrak B}
\newcommand{\Id}{\mathrm{Id}}
\newcommand{\HS}{\mathcal H}
\newcommand{\sigS}{\sigma_S}
\newcommand{\rhoS}{\rho_S}
\DeclareMathOperator{\ran}{ran}
\DeclareMathOperator{\dist}{dist}
\DeclareMathOperator{\supp}{supp}
\DeclareMathOperator{\Rea}{Re}
\DeclareMathOperator{\Ima}{Im}
\title[Weyl's Theorem for Bounded Self-Adjoint Operators]{Weyl's Theorem for Bounded Self-Adjoint Operators: the restriction lemma for spectral projections (printed as Lemma 4.3)}
\author{Hatem Baloudi, Mohamed Ali Dbeibia, Kamel Mahfoudhi}
\date{Source: arXiv:2609.23682v1 (CC BY 4.0)}
\begin{document}
\maketitle

\noindent\emph{Source and licence.} Weyl's Theorem for Bounded Self-Adjoint Operators. Version arXiv:2609.23682v1 (CC BY 4.0), distributed under the Creative Commons Attribution 4.0 International licence, http://creativecommons.org/licenses/by/4.0/, as recorded in the arXiv licence metadata for this exact version and shown on the saved abstract page https://arxiv.org/abs/2609.23682v1. Full licence notice: licenses/arxiv-2609.23682-CC-BY-4.0.txt.\par\smallskip

\noindent\textit{Editorial scope.} Counted unit: the restriction lemma for spectral projections, printed in the article as Lemma 4.3 (source0.tex lines 595--608, one \texttt{lemma} environment with source label \texttt{lem:restriction}) and described by the article itself as the technical core of the paper, delivered together with the article's complete original proof of it (source0.tex lines 610--681, one \texttt{proof} environment of 72 lines immediately adjacent to the statement, with no gap and no deferral). The proof is the article's own argument in the quaternionic spectral theory setting: it uses the spectral measure and the scalar measures attached to the vectors to determine the support of the measure on a compact open part of the spectrum, identifies the range and the kernel of the corresponding spectral projection, proves the closedness, orthogonality and invariance of the two pieces and the direct-sum decomposition, and computes the spectra of the two restrictions together with the positive distance between them in the nonempty case. No mathematics is written, repaired or re-derived: statement and proof are the article's own text line for line, in the article's own notation, and no line of either is omitted, substituted or re-flowed.

Required context retained uncounted: the article's spectral theorem with its labelled display defining the scalar measures (source0.tex lines 560--574), which the delivered proof invokes. Nothing else from the article is needed.

References. The retained context theorem is stated in the article with a citation in its own header, and the three references it names are exactly the citation commands that occur in the retained lines. The package therefore carries exactly those three bibliography entries, transcribed from the article's own bibliography and printed with the article's own labels 2, 11 and 21.

Editorial formatting only, in addition: an AMS \texttt{amsart} document that restates verbatim the article's own macro definitions needed by the retained lines and the article's own theorem-environment declarations with the same shared counter. The article's own class is the standard amsart class, so no publisher class file is involved. Consequently the retained environments print here in the order in which they occur, so the counted lemma prints as Lemma 2 whereas the article prints it as Lemma 4.3 under its per-section numbering, and the retained spectral theorem prints as Theorem 1; the equations of the retained spans are numbered anew in order of appearance and every cross-reference inside the retained text resolves to this wrapper's numbering. That is the only change to the numbering. No figure, no image-inclusion line and no drawing code occur in any retained line, so no artwork is delivered or needed.

Not retained: the article's abstract, its introduction, its preliminaries, its other lemmas, propositions and theorems, its applications and its bibliography except the three cited entries; none of that material is used as a step of the delivered proof, and the delivered proof is complete as the author wrote it. The package ships the article's concatenated source verbatim as \texttt{sources/211134/source0.tex}, and every delivered excerpt is an exact line range of that file. No mathematical statement, hypothesis, estimate or conclusion has been rewritten, repaired or added, and printed slips, if any, are preserved literally.
\begin{theorem}[{Spectral theorem; \cite[Corollary~11.2.2(ii)]{ColomboGantnerKimsey2018},
see also \cite{GhiloniMorettiPerotti2013,AlpayColomboKimsey2016}}]
Let $T=T^*\in\BB(\HS)$. There is a unique projection-valued measure
$E:\Bor(\RR)\to\BB(\HS)$, with values orthogonal projections, such
that $$T=\int_\RR t\,dE(t),$$ it satisfies $\supp(E)=\sigS(T)$. For
every bounded Borel $f:\RR\to\RR$ the operator $$f(T)=\int_\RR f\,dE$$
is bounded self-adjoint, $f\mapsto f(T)$ is a unital homomorphism of
real algebras, and
\begin{equation}
 \|f(T)x\|^2=\int_\RR|f|^2\,d\mu_x,
 \qquad
 \mu_x(B)=\langle E(B)x,x\rangle=\|E(B)x\|^2 .
 \label{eq:scalar}
\end{equation}
\end{theorem}

\begin{lemma}
\label{lem:restriction}
Let $\Delta\subset\sigS(T)$ be nonempty, compact and open in
$\sigS(T)$, and put $$\Delta'=\sigS(T)\setminus\Delta,
\HS_1=\ran(E(\Delta)), \HS_2=\ker(E(\Delta)).$$ Then $\HS_1,\HS_2$
are closed, mutually orthogonal, $T$-invariant,
$\HS=\HS_1\oplus\HS_2$, and
\[
 \sigS(T|_{\HS_1})=\Delta,\qquad \sigS(T|_{\HS_2})=\Delta' .
\]
Moreover $\HS_1\neq\{0\}$, and
$\dist(\sigS(T|_{\HS_1}),\sigS(T|_{\HS_2}))>0$ when
$\Delta'\neq\varnothing$.
\end{lemma}

\begin{proof}
The operator \(E(\Delta)\) is an orthogonal projection commuting with \(T\). Hence
\[
\HS_1=\ran E(\Delta),\qquad \HS_2=\ran E(\Delta')
\]
are closed, orthogonal, \(T\)-invariant, and \(\HS=\HS_1\oplus\HS_2\). Since \(\Delta\) is compact and closed in \(\sigS(T)\), its complement \(\Delta'\) is compact and open in \(\sigS(T)\). Also \(E\) is carried by $$\supp(E)=\sigS(T),$$ so
\[
E(\RR\setminus\Delta)=E(\Delta'),
\]
and therefore \(\HS_2=\ran(E(\Delta'))\). For \(q\in\CC\), define
\[
\varphi_q(t)=t^2-2\Rea(q)t+|q|^2
=\bigl(t-\Rea(q)\bigr)^2+|\Ima(q)|^2 .
\]
Then \(Q_q(T)=\varphi_q(T)\) in the Borel functional calculus. We treat \(\Delta\), the argument for \(\Delta'\) is identical. If \(\Delta'=\varnothing\), then \(\HS_2=\{0\}\) and \(\sigS(T|_{\HS_2})=\varnothing\).

\emph{Step 1: \(q\notin\Delta\) implies \(q\in\rhoS(T|_{\HS_1})\).}
If \(q\notin\RR\), then \(\varphi_q\ge |\Ima(q)|^2>0\). If \(q\in\RR\) and \(q\notin\Delta\), then
\[
\varphi_q(t)=(t-q)^2\ge \dist(q,\Delta)^2>0
\]
on the compact set \(\Delta\). Thus in either case
\[
\delta:=\inf_{\Delta}\varphi_q>0.
\]
Hence \(\psi:=\varphi_q^{-1}\chi_\Delta\) is bounded Borel with \(\|\psi\|_\infty\le\delta^{-1}\), and
\[
\varphi_q\psi=\chi_\Delta
\]
gives
\[
Q_q(T)\,\psi(T)=\psi(T)\,Q_q(T)=E(\Delta).
\]
Both operators preserve \(\HS_1\), and \(E(\Delta)|_{\HS_1}=\Id_{\HS_1}\). Therefore \(\psi(T)|_{\HS_1}\) inverts \(Q_q(T|_{\HS_1})\), so \(q\in\rhoS(T|_{\HS_1})\).

\emph{Step 2: \(q\in\Delta\) implies \(q\in\sigS(T|_{\HS_1})\).}
Such a \(q\) is real. Let \(\varepsilon>0\) and
\[
B_\varepsilon=(q-\varepsilon,q+\varepsilon)\cap\Delta .
\]
Since \(\Delta\) is only relatively open in \(\sigS(T)\), write \(\Delta=W\cap\sigS(T)\) with \(W\subset\RR\) open. Then
\[
O=(q-\varepsilon,q+\varepsilon)\cap W
\]
is open in \(\RR\) and contains \(q\in\supp(E)\). Hence \(E(O)\neq0\). Moreover,
\[
E(O)=E(O\cap\sigS(T))=E(B_\varepsilon),
\]
because \(E\) is carried by \(\supp(E)\). Choose a unit vector
\[
x\in\ran(E(B_\varepsilon))\subset\HS_1 .
\]
Then \(\mu_x\) is carried by \(B_\varepsilon\), and on \(B_\varepsilon\),
\[
|\varphi_q(t)|=(t-q)^2<\varepsilon^2 .
\]
Thus
\[
\|Q_q(T|_{\HS_1})x\|^2\le\varepsilon^4
\]
by \eqref{eq:scalar}. Since \(\varepsilon>0\) is arbitrary, \(Q_q(T|_{\HS_1})\) is not bounded below. Hence \(q\in\sigS(T|_{\HS_1})\).

Steps 1 and 2 yield
\[
\sigS(T|_{\HS_1})=\Delta .
\]
Applying the same argument to \(\Delta'\) gives
\[
\sigS(T|_{\HS_2})=\Delta' .
\]
Since \(\Delta\neq\varnothing\), Step 2 supplies a unit vector in \(\HS_1\), so \(\HS_1\neq\{0\}\). Finally, \(\Delta\) and \(\Delta'\) are disjoint compact sets.
\end{proof}

\begin{thebibliography}{99}
\bibitem[2]{AlpayColomboKimsey2016} D. Alpay, F. Colombo and D. P. Kimsey, The spectral theorem for quaternionic unbounded normal operators based on the $S$-spectrum, \emph{J. Math. Phys.} \textbf{57} (2016), no. 2, 023503.
\bibitem[11]{ColomboGantnerKimsey2018} F. Colombo, J. Gantner and D. P. Kimsey, \emph{Spectral Theory on the $S$-Spectrum for Quaternionic Operators}, Operator Theory: Advances and Applications, vol. 270, Birkh\"auser/Springer, Cham, 2018.
\bibitem[21]{GhiloniMorettiPerotti2013} R. Ghiloni, V. Moretti and A. Perotti, Continuous slice functional calculus in quaternionic Hilbert spaces, \emph{Rev. Math. Phys.} \textbf{25} (2013), no. 4, 1350006.
\end{thebibliography}
\end{document}
